\subsection{Classification}
\label{sec:classify}

We classify each injection CVE along two axes: the \emph{weakness family}
(rooted in the CWE taxonomy) and the injection \emph{syntax type}: the
embedded language the attack targets, such as SQL, HTML, shell, or LDAP.
The syntax axis is the one that matters for matchertext, because containment
depends on how the untrusted value is delimited in the target language, not on
how the weakness is categorized.

\paragraph{Injection gate}
A CVE is injection-related if an injection phrase rule matches its description
or it has a CWE in the \mbox{CWE-74} (``Injection'') subtree, from the CVE record
and NVD, else from other sources; the subtree excludes \mbox{CWE-99} (the value
names a resource) and descendants, and \mbox{CWE-94} is dropped where the description names memory
corruption.
This yields \textbf{\dataInjectionCVEs} injection CVEs, \textbf{\dataInjectionSharePrecise} of published
records, with \dataInjectionSubtreeCWEs CWEs of the subtree (\dataInjectionCWEs with
those co-assigned).
Its neighbors, and the OWASP CWEs outside it, name no language a value is
parsed in (\cref{sec:appendix:corpus}).

\paragraph{Layered labeling}
Each CVE's syntax type is resolved by a fallback chain, ordered from most to
least reliable.
This mirrors the content-classifier methodology of the matchertext empirical
study~\cite{matchertext}.
First, a deterministic map from specific CWEs (\eg \mbox{CWE-89}$\rightarrow$SQL,
\mbox{CWE-79}$\rightarrow$HTML); where one source gives several, the CWE whose
syntax the description names wins, else the lowest CWE number.
Second, phrase rules over the CNA description (``SQL injection'', \dots); the
phrase named first wins.
Third, a word naive-Bayes classifier, trained on the labeled set and on the
gate's rejects as a not-injection class, labels the rest above a margin.
Fourth, an explicit \texttt{unknown} bucket.
Two \emph{delivery rules} refine the label: LLM prompt (A4) where a language
model writes the injected syntax, HTML document (A2) for an XSS whole-document
upload.
Of the injection CVEs, \dataMethodCwe (\dataMethodCweShare) are labeled deterministically from CWEs,
\dataMethodRule from phrase rules, \dataMethodNb from naive Bayes, and only \dataMethodUnknown remain
unknown: the deterministic layers label \dataMethodDeterministicShare.

\paragraph{Label accuracy}
Two raters labeled a blind sample of \dataAuditPerMethod CVEs per layer,
unknown included, from the description alone (Cohen's
$\kappa=\dataAuditKappa$; \dataAuditGraded graded, the rest unclear).
The exact syntax is right for \dataAuditPrecisionCwe of CWE-map labels,
\dataAuditPrecisionRule of rule labels, and \dataAuditPrecisionNb of naive-Bayes
labels, \dataMethodNbShare of the corpus (95\% intervals).
Weighted to the corpus, \dataAuditHostableShare of the CVEs that the raters judge
injection target a matcher-delimitable syntax, above the labeled share of
\cref{sec:prevent}; resolving the \dataLabelConflictCVEs CVEs whose CWEs span the boundary away from
it lowers that share from \dataHostableShareExact to \dataHostableLowShare.

\paragraph{Weakness families}
Aggregated to families, the corpus is dominated by rendering injection
(primarily cross-site scripting) and query injection (majority of SQL),
followed by command injection and code/expression injection.
Because family and syntax are distinct axes (CRLF header injection and XSS
share the rendering family but differ sharply in matchertext-delimitability), the
containment analysis works from the finer syntax axis throughout.
